% !TeX root = ../机器学习笔记.tex
% 由主文档通过 \input 载入；请编译主文档。
\setcounter{chapter}{24} % 保留原书章号。
\chapter{综合实验与可靠机器学习实践}
\label{chap:ml-practice}

完整实验应把任务、数据、方法与证据连起来。下面给出可实施的实验协议及其理论基线，实验数值须由实际运行得到，不能把模型预期表现当作观测结果。

\section{从任务描述到可检验问题}

首先写清预测时刻、输入可用性、输出对象、采样单位和错误代价。例如预测用户未来一周是否点击，与解释过去点击记录，是不同任务；即使使用同一数据表，其特征时效与划分规则也不同。再确定主要评价指标以及可接受的训练、存储和推断成本。

先用常数预测或简单规则检查流程。复杂模型必须证明相对于这个基线增加了什么价值，而不只是达到一个孤立分数。模型改善的是观测分布下的预测，若要推断干预某项因素的效果，需要额外的因果设计，不能从特征重要性直接推出。

\section{实验流程、复现与消融}

每次实验记录数据版本、样本划分、预处理拟合范围、超参数候选、训练停止条件、随机种子和评价脚本约定。重复随机训练刻画算法波动，重新抽样数据刻画样本波动；二者应分别记录。方法比较共享评价样本时，可分析逐样本损失差，重采样仍须尊重相关分组。

消融实验（ablation study）移除或替换一个成分，其他条件尽量保持一致。若移除成分后仍需调参，应对两种方法提供可比较的选择预算，而非只精调新方法。保留失败和异常运行的原因，预先规定哪些运行可排除，避免根据测试结果选择性报告。

\section{回归与分类综合项目}

回归项目可独立生成 $X_1,X_2,Z\sim N(0,1)$，令 $X_3=X_1+0.1Z$，$\varepsilon\sim N(0,0.25)$ 独立于输入，并设
\begin{equation}
 Y=2X_1-X_2+\tfrac12(X_2^2-1)+\varepsilon.
 \label{eq:ml-regression-project}
\end{equation}
一次实验生成一千个样本，固定分为训练六百、验证二百、测试二百，所有特征变换仅在训练中拟合。比较训练均值、线性回归、岭回归、含 $X_2^2$ 的线性模型与树集成；先在验证部分选择配置，再一次性评价测试均方误差。

该生成机制有解析参照：真实条件均值已知，Bayes 均方误差为 $0.25$；最优总体线性预测为 $2X_1-X_2$，其均方误差为 $0.75$，因为中心化平方项与线性特征正交且方差为 $0.5$；最佳总体常数为零，风险为 $5.75$。这些是总体理论值，有限样本拟合和测试结果会波动，不能要求每次测试严格按该顺序排列。

分类项目令 $X_1,X_2$ 独立标准正态，$Y\mid X\sim\operatorname{Bernoulli}(\sigma(-2+1.5X_1-X_2))$。比较多数类、逻辑回归、近邻和树集成，报告对数损失、精确率--召回率表现及指定错误代价。若漏判代价为四、误报为一，真实校准后验的理论阈值为 $0.2$；估计模型的实际阈值选择仍要使用验证数据。

\section{无监督与表示学习综合项目}

聚类项目可从三成分等权二维 Gaussian 混合采样，均值为 $(0,0),(3,0),(0,3)$，协方差先取 $0.2I$，再改为不同方向的伸长椭圆。比较 $K$ 均值、Gaussian 混合与密度聚类，统一输入表示，记录内部指标与基于生成成分的外部指标。由于成分可能重叠，生成标签也不一定能由输入完全恢复。

降维项目从低维信号经线性映射并加噪得到高维输入，检查主子空间、重构误差和下游预测。表示学习项目使用固定训练验证测试集，对比原始特征、随机特征与训练表示；冻结或微调的规则、预测头规模和标签预算必须一致。低维图像可辅助观察，但不能替代定量评价和泄漏检查。

\section{解释、不确定性与责任边界}

全局解释描述总体规律，局部解释说明单个预测；两者都受模型和输入分布限制。置换重要性衡量依赖关系，不能单独证明因果效应。置信分数需要通过校准与实际错误检验，不确定性估计还应在分布变化情形下检查。

分组评价可揭示整体指标掩盖的差异，但不同公平性指标表达不同目标。例如要求各组真正例率相等与要求概率校准相同，在某些分布下可能不能同时满足。应说明采用的目标与依据，避免将一个统计指标当作完整伦理判断。数据授权、个人信息使用与适用规则需要在实际项目中另行核对，本书不以数学指标替代这些要求。

\section{部署、监测与最终学习检查}

部署前固定从原始输入到预测输出的整条流程，检查训练与服务中的词表、缺失处理、尺度、模型状态和输出阈值一致。时延与资源测量需注明硬件、批量和输入规模；离线平均分无法说明所有请求的尾部延迟，也不能直接预测上线收益。

运行后监测输入质量、预测分布、实际有标签效果及关键分组表现。输入变化不等于效果一定下降，标签到达较慢时也不能立即确认模型仍然有效；重新训练和替换模型应使用保留的评价协议。完成全书学习的标准，是能够为一个新任务陈述假设、选择基线、推导核心更新、设计可信评价并解释失败原因，而不只是调用更多算法。



